\section{Discussion}
\label{sec:discussion}

We organise the discussion around the three negative findings, in
descending order of how much they affect what other practitioners should
do.

\subsection{Generator-Family Identity Dominates Evasion}
\label{ssec:family_disc}

Across one detector architecture, one training corpus, and one held-out
test set, evasion ranged from \(0.007\) (BioMistral fakes) to
\(0.987\) (SeqGAN fakes) to \(0.983\) (RAID's mixed LLMs). Three
points follow.

\paragraph{The effect dwarfs everything else we measured.}
Per-round evasion variation in Condition~B (\(0.993, 0.980, 0.987\))
sits inside the \(\pm 0.062\) Wilson CI on \(n=250\); the
between-condition retraining effect is therefore zero within noise. The
agent-rewrite effect (\(\Delta_{\mathrm{evasion}} \approx 0.006\) between
C-round-1 and C-round-3) is similarly inside CI. The
between-generator-family effect (\(\Delta_{\mathrm{evasion}}
\approx 0.98\)) is two orders of magnitude larger than either. If a
practitioner is choosing where to invest engineering time, the lesson is
unambiguous: \emph{the choice of fakes the detector trains against
matters far more than how many adversarial rounds you run, or whether
you wrap an LLM agent around the loop}.

\paragraph{Why SeqGAN is so much harder.}
SeqGAN produces text whose token-level distribution is far from the
biomedical-English manifold the detector was pretrained on. This is
visible in Condition~B's metric pattern: AUC remains acceptable
(\(\approx 0.95\)), meaning the detector \emph{ranks} SeqGAN fakes
above PubMed reals, but F1 collapses to \(0.65\), meaning the
calibrated threshold lands in a region where most SeqGAN samples are
called real. Two epochs of retraining per round do not move the
threshold enough to recover, and from inspection of the per-round
checkpoints they do not move the underlying decision surface either.
This is consistent with the family-of-generators observation in
Grover \cite{zellers2019defending}: the best detector for an
out-of-family generator is often a detector that was \emph{also}
trained on that family.

\paragraph{Disclosure recommendation.}
We argue that any deployed biomedical AI-text detector should publish a
generator-family evasion matrix as part of its documentation, with at
least one in-family LLM, one out-of-family LLM, and one non-LLM
generator. Reporting only in-family AUC produces numbers that look
near-perfect and are operationally meaningless if the deployment
adversary uses any other generator.

\subsection{The C\,$=$\,D Pathology and Best-Val-AUC Selection}
\label{ssec:cd_path}

Conditions~C and~D produced byte-identical test metrics across all three
rounds. We diagnose this as an interaction between the detector's
ability and the silent default of best-validation-AUC checkpoint
selection.

The Condition~A baseline detector reaches validation AUC \(0.9990\) —
within \(\sim 0.001\) of the empirical ceiling on this corpus. In
Condition~D, each retraining round computes the validation AUC of
every epoch checkpoint and keeps the best. Because no retrained
checkpoint can exceed the baseline's \(0.9990\), the selector keeps
rolling back to the round-0 weights — i.e., to Condition~A's
detector. Condition~D's loop \emph{runs}, generates fakes, computes
hard subsets, calls the agent, and writes \(936\) new fakes to disk;
none of that work updates the deployed detector.

This is not a bug in our implementation. It is the expected behaviour
of best-val-AUC selection when the baseline is at ceiling. The
implication is that the metric one reports for ``the full pipeline''
is actually the metric of a detector that no rewrite or retrain ever
modified — an effect that, to our knowledge, is not flagged in the
adversarial-text-training literature.

\paragraph{Two ways out.}
\begin{enumerate}[leftmargin=*]
\item \emph{Adversarial-aware checkpoint selection}: keep the checkpoint
  with the lowest evasion on a held-out adversarial pool, not the
  highest validation AUC. This trades a small amount of clean-data
  performance for measurable adversarial signal.
\item \emph{Stronger generators}: an LLM that has been LoRA-fine-tuned
  to actively evade the detector (e.g., DPO with detector logits as
  reward) would push validation AUC below the baseline ceiling and
  give retraining a non-zero gradient. Our flat
  Condition~B curve suggests that without one of these two changes,
  the loop has nothing to learn from at this scale.
\end{enumerate}

\subsection{Cross-Domain Brittleness}
\label{ssec:domain_disc}

RAID's biomedical slice exposes a calibration failure that the
in-distribution metrics cleanly hide. AUC stays at \(0.82\), so the
detector retains ranking ability across domains. Accuracy collapses to
\(0.11\) and evasion shoots to \(0.983\), so its calibrated threshold
is meaningless on RAID's mixed-LLM fakes. The qualitative pattern is
the same brittleness reported by the RAID and MAGE papers; the
quantitative shape (huge AUC--accuracy gap) is a useful signal for
practitioners who only ever look at AUC.

For the BioShield deployment story, this means we cannot honestly
position the detector as a universal AI-text detector. We position it
instead as a \emph{biomedical specialist} trained against the
BioMistral family, with explicit cross-domain failure documented up
front. Reviewers asking for cross-domain robustness should be pointed
to the multi-family training proposal in Section~\ref{sec:limitations}.

\subsection{Retraction of the 500-Scale ``Adversarial Training Works''
Curve}
\label{ssec:retraction}

A previous version of this work, run at \(500\)-sample scale
(\(250\) real \(+\) \(250\) BioMistral fake, fake-pool
\(n = 250\) per round, test \(n = 100\)), reported a per-round
evasion curve of \(0.74 \to 0.64 \to 0.54\) on Condition~B and
interpreted it as evidence that 3-round adversarial retraining works
against SeqGAN-style fakes. The narrative slide deck for that run
opened with this number.

At \(3{,}000\)-sample scale, with the same
random seed, hyperparameters, and pipeline code, that curve does not
reproduce. Per-round evasion is \(0.993, 0.980, 0.987\) — flat within
the \(\pm 0.062\) Wilson CI on \(n = 250\). We retract the prior
``adversarial training works'' claim and attribute the \(500\)-scale
trend to small-sample noise: the per-round evasion sample is binary on
\(n = 250\), giving an effective standard deviation of
\(\sigma \approx 0.03\), which can produce a \(\Delta = 0.20\) drift
across three rounds with non-trivial probability under the null.

This is, in our reading, the most important methodological lesson of
the project: \emph{adversarial training results at sub-thousand sample
scale should be treated as exploratory, not load-bearing}, and any
published evasion curve should disclose its sample size and the CI
implied by it.
